% TOP_PROBLEM: {"problemId": "problem.termination-of-flips-conjecture-in-the-minimal-model-program", "canonicalTitle": "Termination of Flips Conjecture in the Minimal Model Program", "releaseRank": 72, "primaryDomain": "geometry_topology", "primaryDomainLabel": "Geometry & topology", "displayStatus": "open_with_solved_subcases", "releaseStatus": "open_with_solved_subcases"}
% !TeX program = lualatex
% Definition and sources only; this file is not an LLM solution attempt.
\documentclass[11pt]{article}
\usepackage[margin=25mm]{geometry}
\usepackage{fontspec}
\setmainfont{Latin Modern Roman}
\usepackage{amsmath,amssymb,mathtools}
\usepackage{unicode-math}
\setmathfont{Latin Modern Math}
\usepackage{xurl,hyperref}
\hypersetup{colorlinks=true,urlcolor=blue}
\setcounter{secnumdepth}{0}
\setlength{\parindent}{0pt}
\setlength{\parskip}{5pt}
\setlength{\emergencystretch}{3em}
\begin{document}
\section*{\texorpdfstring{Termination of Flips Conjecture in the Minimal Model Program}{Termination of Flips Conjecture in the Minimal Model Program}}\label{72}
\subsection{Definitions and mathematical statement}
A log pair $(X,\Delta)$ has $X$ normal, $\Delta$ an effective boundary divisor, and $K_X+\Delta$ real or rational Cartier in the chosen MMP category. Log canonicity requires all discrepancies to be at least $-1$; klt requires them to be greater than $-1$. A flip replaces a small negative contraction by the small birational model on which the transformed adjoint divisor is relatively ample. The claim excludes any infinite chain
\[
 (X_0,\Delta_0)\dashrightarrow(X_1,\Delta_1)\dashrightarrow\cdots
\]
of the allowed adjoint flips starting from a projective complex log canonical pair. A theorem about one prescribed choice of flips is not automatically termination of every sequence.
\par\smallskip\textit{Formulation and concept references:} \href{https://master-math-fonda.imj-prg.fr/gt/mmp.pdf}{[S1]}.

\subsection{Short English statement}
Must every allowed sequence of flips in the complex minimal model program stop after finitely many steps?
\subsection{Sources}
\begin{itemize}
\item[S1]Vladimir Lazic, Zhixin Xie, "Nakayama-Zariski decomposition and the termination of flips," arXiv:2305.01752v4; received July 8, 2024, accepted June 25, 2025.
\url{https://master-math-fonda.imj-prg.fr/gt/mmp.pdf}.
\textit{Formulation source.}
\item[S2]MMP for Enriques pairs and singular Enriques varieties.
\url{https://jep.centre-mersenne.org/item/10.5802/jep.334.pdf}.
\textit{Further reference; not a proof endorsement.}
\item[S3]Nakayama-Zariski decomposition and the termination of flips.
\url{https://epiga.episciences.org/16998}.
\textit{Further reference; not a proof endorsement.}
\item[S4]On termination of flips and exceptionally non-canonical singularities.
\url{https://arxiv.org/abs/2209.13122}.
\textit{Further reference; not a proof endorsement.}
\end{itemize}
\end{document}
